\section{Intermediate Encoding Format}\label{sec:inter-parse}
The simple parsing format described in Section~\ref{sec:simpl-parse} was
designed to show that we could reason about nested messages. However, there are
more complex parts of Protobuf we need to support before being able to verify
the whole thing. Some of these features were introduced to an intermediate
complexity encoding format, which we also implemented and verified.

These important features were added to the intermediate encoding format:
\begin{itemize}
\item Variable number of fields per message.
\item Encoding uses key-value pairs which can be serialized in any order.
\item Drops unknown fields without issue.
\item Does not encode, but can recover, missing / default values.
\item Parsing requires a descriptor to differentiate values.
\end{itemize}

\subsection{Type Definitions}\label{sec:ip-types}
Below is the Lean definitions for a descriptor and message. Unlike the original
Rocq code, which use a \texttt{gmap} from the \texttt{stdpp}, Lean's finite maps
cannot be nested into an inductive type. Instead, the intermediate format models
the descriptors and a list of key-value pairs, then requires theorems over them
maintain a well-formedness invariant that the list is sorted and contains no
duplicates. Notice that these are mutually inductive types, with a
\texttt{Field} able to hold a \texttt{Desc} for a nested message.

\begin{minted}{lean}
-- A descriptor maps field numbers (integers) to field types.
-- Mutually defined with `Field` which may contain nested descriptors.
--- Corresponds to `Desc`/`Field` from `InterParse/Descriptor.v`.
mutual
inductive Desc where
  | mk (fs : List (Int × Field))
inductive Field where
  | msg (d : Desc)
  | bool
  | int
end

-- A value maps field numbers to typed values.
-- Mutually defined with `Val` which may contain nested values.
-- Corresponds to `Value`/`Val` from `InterParse/Descriptor.v`.
mutual
inductive Value where
  | mk (vs : List (Int × Val))
inductive Val where
  | msg (v : Value)
  | bool (b : Bool)
  | int (z : Int)
  | missing
end
\end{minted}

Represented mathematically, you could describe these types like so:

\vspace{1em}
\begin{definition}[Intermediate Descriptor]\label{def:ip-desc}
 An intermediate descriptor is a map from integers representing the field number
 to a field definition, which can either be
 \begin{itemize}
 \item A tag denoting a nested message, which also has a nested intermediate descriptor.
 \item A tag denoting a boolean.
 \item A tag denoting an integer.
 \end{itemize}
\end{definition}

\vspace{1em}
\begin{definition}[Intermediate Message]\label{def:ip-msg}
  An intermediate message is a map from integers representing the field number
  to a values, which can either be
  \begin{itemize}
  \item A nested message, without any descriptor.
  \item A boolean.
  \item An integer.
  \item A special missing value.
  \end{itemize}
\end{definition}

The primary difference between descriptors and messages is that messages have
real values in them while descriptors have tags without values and describe a
potential shape of a message. You can thing of descriptors as defining a formal
type of message.

\subsection{Encoding}\label{sec:ip-encoding}

When encoding a message under the intermediate format, each value is encoded
accordingly:
\begin{itemize}
\item \textbf{Booleans:} Encoded into four little-endian bytes with zero for
  false and any positive integer (typically 1) for true.
\item \textbf{Integers:} Encoded into four little-endian bytes. Currently only
  positive integers are supported, but wrapper functions could support negative
  integers by converting them to the equivalent two's-complement positive
  integer.
\item \textbf{Nested Messages:} A length prefixed nested list of key-value
  pairs. The length is encoded as a one byte integer, limiting the format's
  ability to handle nested messages to messages which can be encoded into 255 or
  fewer bytes. (Note that the top-level message can exceed 255 bytes, but that
  couldn't later be put into a different message.)
\end{itemize}

Each field is then encoded first with the one byte field ID number, then the
encoding described above. A message is a list of key-value pairs. Assuming every
field you have encoded is in the descriptor, they can all be parsed in a
self-contained manner.

\vspace{1em}
\begin{definition}[Intermediate Descriptor Well-Formedness]\label{def:inter-dwf}
  A descriptor $\dc$ is considered well-formed, denoted $\wfd{\dc}$, if
  \begin{itemize}
  \item The keys are stored in sorted order.
  \item There are no duplicate keys.
  \end{itemize}
\end{definition}

\vspace{1em}
\begin{definition}[Intermediate Value Well-Formedness]\label{def:inter-vwf}
  A message $\m$ is well-formed, denoted $\wfm{\m}$, if
  \begin{itemize}
  \item All keys are in the range 0 to 256.
  \item All values in the message are type-aligned with the descriptor
    (i.e.\ integer fields in the descriptor have integer values in the message).
  \item All integer payloads are within the range of a valid 32-bit unsigned
    integer.
  \item All nested messages are also well-formed.
  \end{itemize}
\end{definition}

\subsection{Missing Values \& Descriptor Mismatches}\label{sec:ip-missing}
The presence of missing values and a dependency on a descriptor complicates our
definition of what it means to successfully serialize and parse. For example, we
expect that if the message being serialized have a field id \emph{not} in the
descriptor we're serializing under, that field will be omitted. If we then parse
under the same descriptor, the recovered value won't be identical to one we put
into the serializer; it will be missing that field.

Recall that a serializer comes with a well-formed condition, and that the
traditional definition for parser-serializer correctness is this:

\vspace{1em}
\begin{definition}[Parser Correctness]\label{def:base-pc}
  A serializer $ser$ with well-formed condition $wf$ is correct with respect to
  a parser $par$ if
  \[\forall\ x\ enc\ rest,\;wf\ x \rightarrow ser\ x = \mathtt{Success}\ enc \rightarrow par\ (enc \doubleplus rest) =
    \mathtt{Success}\ x\ rest \]
\end{definition}

To account for descriptor mismatches, and to pave the way for a compatibility
respecting parser correctness definitions in the future, this definition can be
expanded to

\vspace{1em}
\begin{definition}[Simple Relational Parser Correctness]\label{def:simp-pc}
  A serializer $ser$ with well-formed condition $wf$ is relational correct with
  respect to a parser $par$ and relation $R$ if
  \[ \forall\ x\ enc\ rest,\;wf\ x \to ser\ x = \mathtt{Success}\ enc \rightarrow \exists\ x',\;par\ (enc
    \doubleplus rest) = \mathtt{Success}\ x'\ rest \wedge R\ x\ x' \]
\end{definition}

Since our format requires a descriptor to differentiate between fields, and does
not support arbitrary suffixes (since the top level message isn't length limited
in any way and you could always have another field encoded in the suffix), we
need a modified version of Definition~\ref{def:simp-pc}.

\vspace{1em}
\begin{definition}[Relational Parser Correctness]
  A serializer $ser$ with well-formed condition $wf$ is relational correct with
  respect to a parser $par$, descriptor $\dc$ and relation $R$ if
  \[ \forall\ x\ enc,\;wf\ x \to ser\ \dc\ x = \mathtt{Success}\ enc \rightarrow \exists\ x',\;par\ \dc\ enc
    = \mathtt{Success}\ x'\ \nil \wedge R\ \dc\ x\ x' \]
\end{definition}\label{def:rpc}

This definition already looks like our definition of compatibility
(Definition~\ref{def:compat}), with just an extra check that the serializer
actually succeeds.

\subsection{Schema-Correct Compatibility Relation}\label{sec:sc-id-compat}

In order to simplify proof burdens for the moment, there is another important
concept here: when a value is \emph{fully described} or \emph{schema correct} by
a descriptor. Intuitively, this means that every value is the message has a
corresponding field in the descriptor and vice versa. This relation is denoted
as $\langle \m : \dc \rangle$ and is defined with these two rules:

\vspace{1em}
\begin{definition}[Schema Correct]
  A value $\m$ is \emph{schema correct} with respect to a descriptor $\dc$ if
  \begin{itemize}
  \item Every value in $\m$ is declared in $\dc$.
  \item Every field in $\dc$ id present in $\m$.
  \item Every value in $\m$ matches the type required by the corresponding field
    in $\dc$.
  \item All nested messages are schema correct with respect to the descriptors
    embedded in $\dc$.
  \end{itemize}
\end{definition}

\begin{mathpar}
  \infer[SC-emp]{ }{\langle \emptyset : \emptyset \rangle}

  \infer[SC-insert-bool]{ \langle \m : \dc \rangle \\ \id \not \in \dom{\m} \\ \id \not \in \dom{\dc}}
  { \langle \m[\id \mapsto \mathtt{M\_BOOL}\ b] : \dc[\id \mapsto \mathtt{F\_BOOL}] \rangle}

  \infer[SC-insert-int]{ \langle \m : \dc \rangle \\ \id \not \in \dom{\m} \\ \id \not \in \dom{\dc}}
  { \langle \m[\id \mapsto \mathtt{M\_INT}\ z] : \dc[\id \mapsto \mathtt{F\_INT}] \rangle}

  \infer[SC-insert-msg]{ \langle \m : \dc \rangle \\ \id \not \in \dom{\m} \\ \id \not \in \dom{\dc} \\ \langle
    \m' : \dc' \rangle}
  { \langle \m[\id \mapsto \mathtt{M\_MSG}\ \m'] : \dc[\id \mapsto \mathtt{F\_MSG}\ \dc'] \rangle}
\end{mathpar}

This relation is designed to ensure that the related message and descriptor are
in lock-step as new fields are added to it.

We can now turn our attention to the compatibility relation built on top of
this. While at the moment we're only interested in having the same descriptor on
both sides of the compatibility relation, it's been written for maximum
flexibility and to generalize to a more complex relation.

\begin{mathpar}
  \infer[Compat-Refl]{\langle \m : \dc_1 \rangle \\ \langle \m : \dc_2 \rangle}{\langle \m : \dc_1 \msgrel \m : \dc_2 \rangle}

  \infer[Compat-Add]{\langle \m_1 : \dc_1 \msgrel \m_2 : \dc_2 \rangle \\\\ \id \not \in \dom{\m_1} \\ \id \not \in
    \dom{\dc_2} \\ \id \not \in \dom{\m_2} \\ \id \not \in \dom{\dc_2} \\\\ v_1 = v_2 \\ \fd_1
    = \fd_2}
  { \langle \m_1[\id \mapsto v_1] : \dc[\id \mapsto \fd_1] \msgrel
    \m_2[\id \mapsto v_2] : \dc_2[\id \mapsto \fd_2] \rangle}
\end{mathpar}

That is also why the \textsc{Compat-Add} rule has premises $v_1 = v_2$ and $\fd_1
= \fd_2$ for a more complex compatibility story, this was call out to the value
relation (Section~\ref{sec:val-rel}) and type relation (Section~\ref{sec:typ-rel}).

\vspace{1em}
\begin{lemma}[Intermediate Schema Correct Equal]\label{lem:inter-sc-eq}
  \[
    \forall\ \m_1, \m_2, \dc_1, \dc_2.\ \langle \m_1 : \dc_1 \msgrel \m_2 : \dc_2 \rangle
    \rightarrow \dc_1 = \dc_2 \rightarrow v_1 = v_2
  \] 
\end{lemma}

\vspace{1em}
\begin{theorem}[Intermediate Schema Correct Round-trip]\label{thm:inter-schema-correct}
  \[
    \langle \m_1 : \dc \rangle \rightarrow \wfm{\m_1} \rightarrow \forall\ \bs, \mathrm{ser}_{\dc}(\m_1) =
    \bs \rightarrow \exists\ \m_2, \mathrm{par}_{\dc}(\bs) = \m_2 \wedge \langle \m_1
    : \dc \msgrel \m_2 : \dc \rangle
  \]
\end{theorem}

The immediate corollary of Lemma~\ref{lem:inter-sc-eq} and
Theorem~\ref{thm:inter-schema-correct} is that the $\m_1$ and $\m_2$ in the
theorem statement are equal.  

\subsection{The Identity Compatibility Relation}\label{sec:id-compat}
Naturally then, we can pose the descriptor mismatch problem in terms of a
compatibility relation where the descriptor doesn't change. Since the descriptor
is always related to itself, this is known as the \emph{identity compatibility
  relation}. Let $\langle \m_1 : \dc_1 \msgrel \m_2 : \dc_2 \rangle$ denote that $\m_1$
when interpreted under $\dc_1$ is compatible with $\m_2$ under $\dc_2$.


\subsubsection{Full Identity Compatibility Relation}\label{sec:full-id-compat}
If we want to relax the requirement that the value is always fulled described by
the paired descriptor, we have to allow for values to be dropped during
serialization and missing values inserted during parsing. In this setting, it
\emph{should} be the case that any field in both the value and descriptor have
matching types.

\begin{mathpar}
  \infer[Empty]{ }{\langle \emptyset : \emptyset \msgrel \emptyset : \emptyset \rangle}

  \infer[Insert-Int]{\dc' = \dc[\id \mapsto \mathtt{F\_INT}] \\ \langle \m' : \dc \msgrel \m : \dc \rangle}{\langle \m'[\id \mapsto z]
    : \dc' \msgrel \m[\id \mapsto z] : \dc' \rangle}

  \infer[Insert-Bool]{\dc' = \dc[\id \mapsto \mathtt{F\_BOOL}] \\ \langle \m' : \dc \msgrel \m : \dc \rangle}{\langle \m'[\id \mapsto b]
    : \dc' \msgrel \m[\id \mapsto b] : \dc' \rangle}

  \infer[Insert-Msg]{\dc' = \dc[\id \mapsto \mathtt{F\_MSG}\ \dc_n] \\\\ \langle \m' : \dc \msgrel \m : \dc \rangle \\ \langle
    \m_n' : \dc_n \msgrel \m_n : \dc_n \rangle}{\langle \m'[\id \mapsto \mathtt{M\_MSG}\ \m_n'] : \dc' \msgrel \m[\id \mapsto
    \mathtt{M\_MSG}\ \m_n] : \dc' \rangle}

  \infer[Drop]{\id \not \in \dom{\dc} \\ \id \not \in \dom{\m} \\ \langle \m' : \dc \msgrel \m : \dc \rangle}{\langle \m'[\id \mapsto z] : \dc \msgrel \m : \dc \rangle}

  \infer[Add-Missing]{ \id \not \in \dom{\m} \\ \dc' = \dc[\id \mapsto \fd] \\ \langle \m' : \dc \msgrel \m :
    \dc \rangle}{ \langle \m' : \dc' \msgrel \m[\id \mapsto \mathtt{M\_MISSING}] : \dc' \rangle}

  \infer[Input-Missing]{\id \not \in \dom{\dc} \\ \id \not \in \dom{\m} \\ \id \not \in
    \dom{\m'} \\\\ \dc' = \dc[\id \mapsto \fd]}{\langle \m'[\id \mapsto \mathtt{M\_MISSING}] : \dc' \msgrel \m[\id \mapsto \mathtt{M\_MISSING}] : \dc' \rangle}
\end{mathpar}

\begin{theorem}[Intermediate Identity Compatibility]\label{thm:inter-id-compat}
  \[
    \wfd{\dc} \rightarrow \wfm{\m_1} \rightarrow \forall\ \bs, \mathrm{ser}_{\dc}(\m_1) =
    \bs \rightarrow \exists\ \m_2, \mathrm{par}_{\dc}(\bs) = \m_2 \wedge \langle \m_1
    : \dc \msgrel \m_2 : \dc \rangle
  \] 
\end{theorem}

Notice that this theorem has a much weaker precondition, just that the
descriptor is well-formed rather than that the value is schema-correct and that
the resulting statement uses a weaker compatibility relation which doesn't have
a squeeze lemma like the schema correct compatibility relation did
(Lemma~\ref{lem:inter-sc-eq}).

\subsection{Full Compatibility Relation}\label{sec:ip-compat-rel}

Now we develop the full compatibility relation for the intermediate encoding
format. Like with the protobuf relations, this will be separated into a value
relation, type relation and message relation.
\begin{align*}
  \m &: \text{message} & v &: \text{value} \\
  \dc &: \text{descriptor} & \fd &: \text{field} \\
  \id &: \text{field number}
\end{align*}
Below is a list of the relations about to be defined, along with an English
explanation of how to read them.
\begin{align*}
  \valrel &: \text{Value relation} & v_1 : \ty_1 \valrel v_2 : \ty_2 &\rightarrow ``\text{value } v_1 : \ty_1 \text{ becomes } v_2 :
                                           \ty_2 " \\
  \fieldrel &: \text{Field type relation} & \fd_1 \fieldrel \fd_2 &\rightarrow ``\fd_1 \text{ is type-compatible with }
                                          \fd_2" \\
  \dcompat &: \text{Descriptor type relation} & \dc_1 \dcompat \dc_2 &\rightarrow ``\dc_1 \text{ is
                                                     descriptor-compatible with } \dc_2" \\
  \msgrel &: \text{Message relation} & \m_1 \msgrel \m_2 &\rightarrow ``\text{message } \m_1 \text{ becomes } \m_2"
\end{align*}

\subsubsection{Value Relation}\label{sec:ip-val-rel}
Since booleans and integers are both encoded in a 32 bit little-endian way, we
can convert between them. Message relation compatible messages can also be used
to make updates to nested messages.

\begin{mathpar}
  \infer[V-Bool-Int]{v_2 = \ind{v_1}}{v_1 : \bool \valrel v_2 : \mathtt{int}}

  \infer[V-Int-Bool]{v_2 = \left\{
    {\begin{array}{l@{\quad:\quad}l}
      \false & \mathrm{if}\ v_1 = 0 \\
      \true & \mathrm{otherwise}
    \end{array}}\right.}{v_1 : \mathtt{int} \valrel v_2 : \bool}

  \infer[V-Msg]{\m_1 : \dc_1 \msgrel \m_2 : \dc_2}{\mathtt{M\_MSG}\ \m_1 : \mathtt{F\_MSG}\
    \dc_1 \valrel \mathtt{M\_MSG}\ \m_2 : \mathtt{F\_MSG}\ \dc_2}

  \infer[V-Refl]{ }{v : \fd \valrel v : \fd}

  \infer[V-Trans]{v_1 : \fd_1 \valrel v_2 : \fd_2 \\ v_2 : \fd_2 \valrel v_3 : \fd_3}{v_1 : \fd_1 \valrel v_3 : \fd_3}
\end{mathpar}

\subsubsection{Field Type Relation}\label{sec:ip-f-typ-rel}
Just like with the value relation, we can freely convert between integer and
boolean fields.

\begin{mathpar}
 \infer[F-Bool-Int]{ }{\mathtt{bool} \typrel \mathtt{int}}

 \infer[F-Int-Bool]{ }{\mathtt{int} \typrel \mathtt{bool}}

 \infer[F-Msg]{\dc_1 \dcompat \dc_2}{\mathtt{F\_MSG}\ \dc_1 \typrel \mathtt{F\_MSG}\ \dc_2}

 \infer[F-Refl]{ }{\fd \fieldrel \fd}

 \infer[F-Trans]{\fd_1 \fieldrel \fd_2 \\ \fd_2 \fieldrel \fd_3}{\fd_1 \fieldrel \fd_3}
\end{mathpar}

The more complex thing here is handling the descriptors. We need to be able to
reason about weather an update to those is compatible without considering
values. I'm not sure if this would be better in a new relation, but for the
moment I've decided to split them out.

\subsubsection{Descriptor Type Relation}\label{sec:ip-d-typ-rel}

\begin{mathpar}
  {\color{red}
  \infer[D-Update]{ \dom{\dc_{mand}} \supseteq \dom{\dc_{mand}'} \\ \forall\ \id \in
    \dom{\dc} \cap \dom{\dc'}.\ \dc[\id] \fieldrel \dc'[\id] }{ \dc \dcompat \dc' }
  }

  \infer[D-Emp]{ }{\emptyset \dcompat \emptyset}

  \infer[D-Add]{\id \not \in \dom{\dc}}{\dc \dcompat \dc[\id \mapsto \fd]}

  \infer[D-Chg]{\dc[\id] \fieldrel \fd}{\dc \dcompat \dc[\id \mapsto \fd]}

  \infer[D-Refl]{ }{\dc \dcompat \dc}

  \infer[D-Trans]{\dc_1 \dcompat \dc_2 \\ \dc_2 \dcompat \dc_3}{\dc_1 \dcompat \dc_3}
\end{mathpar}

Just like with the Protobuf relation, we can add fields or change existing
fields. This is like width and depth sub-typing for record types. The
\textsc{D-Update} incorporates both width and depth sub-typing on mandatory
fields (not yet supported by the Pollux framework) and that all fields which
align between the descriptors have a compatible type.

It is worth noting that we \emph{cannot} add a \textsc{D-Drop} rule to this
relation. When parsing a field that's not in the descriptor, there is currently
no why to distinguish between a primitive (integer or Boolean) and a nested
message in the wire format. We'd parse the field ID and realize we need to skip
the field, but not know if the field is a 4 byte primitive or a nested message
of dynamic length. This restriction will be fixed in a full protobuf
implementation.

\subsubsection{Message Relation}\label{sec:ip-msg-rel}

\begin{mathpar}
  \infer[M-Emp]{ }{\emptyset : \emptyset \msgrel \emptyset : \emptyset}

  {\color{red}\infer[M-Add]{\m_1 : \dc_1 \msgrel \m_2 : \dc_2 \\ v : \fd \\\\ \id \not \in \dom{\m_1} \\ \id \not \in
    \dom{\dc_1}}{ \m_1 : \dc_1 \msgrel \m_2[\id \mapsto v] : \dc_2[\id \mapsto \fd]}}

  \infer[M-Missing]{\m_1 : \dc_1 \msgrel \m_2 : \dc_2 \\\\ \id \not \in \dom{\m_1} \\ \id \not \in
    \dom{\dc_1}}{ \m_1 : \dc_1 \msgrel \m_2[\id \mapsto \mathtt{M\_MISSING}] : \dc_2[\id \mapsto \fd]}

  \infer[M-Declare]{\m_1 : \dc_1 \msgrel \m_2 : \dc_2 \\\\ \id \not \in \dom{\dc_1} \\ \id \not \in
    \dom{\m_1} \\ \fd_1 \fieldrel \fd_2}{ \m_1 : \dc_1[\id \mapsto \fd_1] \msgrel \m_2[\id \mapsto \mathtt{M\_MISSING}] :
    \dc_2[\id \mapsto \fd_2]}

  \infer[M-Update]{\m_1 : \dc_1 \msgrel \m_2 : \dc_2 \\\\ \m_1[\id] : \dc_1[\id] \valrel v : \fd \\
    \dc_1[\id] \fieldrel \fd}{ \m_1 : \dc_1 \msgrel \m_2[\id \mapsto v] : \dc_2[\id \mapsto \fd]}

  \infer[M-Drop]{\m_1 : \dc_1 \msgrel \m_2 : \dc_2 \\\\ \id \not \in \dom{\m_2} \\ \id \not
    \in \dom{\dc_2}}{ \m_1[\id \mapsto v] : \dc_1[\id \mapsto \fd] \msgrel \m_2 : \dc_2 }

  \infer[M-Drop-Unknown]{\m_1 : \dc_1 \msgrel \m_2 : \dc_2 \\\\ \id \not \in \dom{\m_2} \\ \id \not
    \in \dom{\dc_2} \\ \id \not \in \dom{\dc_1}}{ \m_1[\id \mapsto v] : \dc_1 \msgrel \m_2 : \dc_2 }

  \infer[M-Refl]{ }{ \m : \dc \msgrel \m : \dc}

  \infer[M-Trans]{\m_1 : \dc_1 \msgrel \m_2 : \dc_2 \\ \m_2 : \dc_2 \msgrel \m_3 : \dc_3}{\m_1 : \dc_1 \msgrel
    \m_3 : \dc_3}
\end{mathpar}

We can add new key-value pairs to a message, or update the value in a message
(assuming we keep the fields in the descriptor updated too). Note that the
premise $v : \fd$ in the rule \textsc{M-Add} denotes that the value $v$ is aligns
with the field $\fd$, i.e. $v$ is an integer and $\fd$ is \verb|F_INT|. It would not
hold if $v = \mathtt{true}$ and $\fd$ was \verb|F_INT|.

The highlighted rule would let the parser pull any correctly typed value out of
thin air. Do we want this? No, the $\msgrel$ relation defines how a message is
transmuted during the serializer / parser round trip, and that's not how the
parser would behave.

In \textsc{M-Declare}, we require $\fd_1 \fieldrel \fd_2$ so that we can handle examples
were an empty message under a descriptor with one boolean parses to an empty
field under a different, compatible type. This continues to mirror how the
parser and serializer would actually behave.

\vspace{1em}
\begin{theorem}[Intermediate Compatibility]\label{thm:inter-compat}
  \[
    \wfd{\dc_1} \rightarrow \wfm{\m_1} \rightarrow \wfd{\dc_2}
    \rightarrow \dc_1 \dcompat \dc_2 \rightarrow \forall\ \bs, \mathrm{ser}_{\dc_1}(\m_1) =
    \bs \rightarrow \exists\ \m_2, \mathrm{par}_{\dc_2}(\bs) = \m_2 \wedge \langle \m_1
    : \dc_1 \msgrel \m_2 : \dc_2 \rangle
  \] 
\end{theorem}

% Local Variables:
% citar-bibliography: ("../pollux.bib")
% TeX-master: "../pollux.tex"
% TeX-command-default: "Make"
% jinx-local-words: "protobuf"
% End:
